% Part: counterfactuals
% Chapter: minimal-change-semantics
% Section: transitivity

\documentclass[../../../include/open-logic-section]{subfiles}

\begin{document}

\olfileid{cnt}{min}{tra}

\olsection{تعدی}

مادی شرطیے لئی لڑی دا قاعدہ قائم اے۔ $!A \lif !B, !B \lif !C \Entails !A \lif !C$۔
دوجے لفظاں وچ مادی شرطیہ متعدی اے۔ کی خلافِ واقعہ شرطیاں لئی وی ایہو گل سچ اے؟
سٹالنیکر دی ہِیٹھلی مثال ویکھو۔
\begin{quote}
  جے جے۔ ایڈگر ہوور روسی جمیا ہوندا، تاں اوہ کمیونسٹ ہو جاندا۔

  جے جے۔ ایڈگر ہوور کمیونسٹ ہوندا، تاں اوہ غدّار ہو جاندا۔

  سو جے جے۔ ایڈگر ہوور روسی جمیا ہوندا، تاں اوہ غدّار ہو جاندا۔
\end{quote}
جے ہوور اوسے ویلے، جدوں اوہ اصل وچ جمیا سی، ریاست ہائے متحدہ دی تھاں روس
وچ جمیا ہوندا، تاں اوہ سوویت یونین وچ وڈا ہوندا تے کمیونسٹ بن جاندا
(ایہ منّ لئیے)۔ سو پہلا مقدمہ سچ اے۔ ایس طرح دوجا مقدمہ وی، جے اوہنوں
کلّا ویکھیے، سچ اے۔ پر نتیجہ جھوٹ اے۔ جے ہوور سوویت یونین وچ جمیا ہوندا،
تاں وڈا امکان ایہ سی کہ اوہ پکا کمیونسٹ ہوندا تے (اپنے ملک دا) غدّار نہ ہوندا۔
سچائی دیاں ایہ قدراں جیہڑیاں سمجھ مطابق دتیاں نیں، سٹالنیکر--لوئس دی تشریح نال
میل کھاندیاں نیں۔ ساڈے جہان دے سبھ توں نیڑلا ممکن جہان، جیہدے وچ اکّو تبدیلی
ہوور دی جنم تھاں اے، اوہ اے جتھے ہوور وڈا ہو کے سوویت یونین دا چنگا شہری
بن جاندا اے۔ ایہ سبھ توں نیڑلا ممکن جہان اے جتھے پہلے مقدمے تے نتیجے دا
مقدم سچ اے۔ اوس جہان وچ ہوور کمیونسٹ پارٹی دا وفادار رکن اے، سو غدّار نہیں۔
پر دوجا مقدمہ پرکھن لئی سانوں اک ہور جہان ویکھنا پینا اے۔ اوہ سبھ توں نیڑلا
جہان جتھے ہوور کمیونسٹ اے، اوہ اے جتھے اوہ ریاست ہائے متحدہ وچ جمیا، پاسہ
بدل لیا تے ایس طرح غدّار بن گیا۔\footnote{مثال دا زور سمجھن لئی سانوں
  وجودیات تے سیاست بارے کجھ مفروضے منّنے پین گے۔ مثلاً ایہ ممکن اے کہ
  ہوور روسی ماپیاں دے گھر جمدا، یا ایہ کہ ویہویں صدی دے پنجاہویں ورھیاں وچ
  ریاست ہائے متحدہ دے کمیونسٹ اپنے ملک دے غدّار سن۔}

\begin{prob}
  خلافِ واقعہ شرطیاں دی تعدی ناکام ہون دی اک منن جوگی، سوکھی سمجھ آؤن والی مثال لبھو۔
\end{prob}

\begin{ex}\ollabel{ex:trans-counterex}
  گولیاں والی معنیات وچ ایہ استنباط غیر معتبر اے، یعنی ساڈے کول $p
  \cif q, q \cif r \Entails/ p \cif r$ اے۔ نمونہ $\mModel{M}
  = \tuple{W, O, V}$ ویکھو، جیہدے وچ $W = \{w, w_1, w_2\}$،
  $O_w = \{\{w\}, \{w, w_1\}, \{w, w_1, w_2\}\}$، $V(p) = \{w_2\}$،
  $V(q) = \{w_1, w_2\}$، تے $V(r) = \{w_1\}$۔ $p$ نوں تھاں دین والا
  گولا $S = \{w, w_1, w_2\}$ موجود اے تے $p \lif q$ اوہدے سارے جہاناں
  اُتے سچ اے، سو $\mSat{M}{p \cif q}[w]$۔ $q$ نوں تھاں دین والا گولا
  $S' = \{w, w_1\}$ وی موجود اے تے $\mSat/{M}{q \lif r}$ اوہدے سارے جہاناں
  اُتے سچ اے، سو $\mSat{M}{q \cif r}[w]$۔ پر $p$ نوں تھاں دین والا گولا
  $\{w, w_1, w_2\}$ وچ اک ایہو جیہا جہان اے، یعنی~$w_2$، جتھے
  $\mSat/{M}{p \lif r}[w_2]$۔
\end{ex}

\begin{prob}
\olref[cnt][min][tra]{ex:trans-counterex} دی ردّی مثال نال ملدا گولیاں دا نقشہ بناؤ۔
\end{prob}

\begin{prob}
  \olref[cnt][min][tra]{ex:trans-counterex} وچ جہان $w_2$ اوہ اے جتھے
  ہوور روس وچ جمیا، کمیونسٹ اے تے غدّار نہیں، تے $w_1$ اوہ جہان اے جتھے
  ہوور ریاست ہائے متحدہ وچ جمیا، کمیونسٹ اے تے غدّار اے۔ ایس نمونے وچ
  $w_1$،~$w$ دے $w_2$ نالوں نیڑے اے۔ کی ایہ ضروری اے؟ کی تسی اک ردّی
  مثال دے سکدے او جیہڑی ایہ نہ منّے کہ ہوور دا روس وچ جمنا اوہدے
  کمیونسٹ ہون نالوں زیادہ دُور دا امکان اے؟
\end{prob}

\end{document}
